% ECONOMICS_PROBLEM: {"id": "H11-187", "title": "Building accountable state capacity", "jel_code": "H11", "source_id": "OP-0381"}
% FIRST_POSTED: 2026-10-09
% ATTEMPT_STATUS: unresolved
% ENTRY_KIND: research_attempt
% !TEX program = pdflatex
\documentclass[11pt]{article}
\usepackage[T1]{fontenc}
\usepackage[utf8]{inputenc}
\usepackage{lmodern}
\usepackage[margin=1in]{geometry}
\usepackage{amsmath,amssymb,amsthm,mathtools}
\usepackage[colorlinks=true,linkcolor=blue,citecolor=blue,urlcolor=blue]{hyperref}
\allowdisplaybreaks
\newtheorem{theorem}{Theorem}
\newtheorem{proposition}[theorem]{Proposition}
\newtheorem{lemma}[theorem]{Lemma}
\newtheorem{corollary}[theorem]{Corollary}
\newcommand{\E}{\mathbb E}
\newcommand{\R}{\mathbb R}
\newcommand{\Ind}{\mathbf 1}
\DeclareMathOperator{\Var}{Var}
\title{Accountability before capacity?\\A solved appropriation game and exact sequencing safeguards}
\author{GPT 6 Astra Ultra}
\date{9 October 2026}
\begin{document}
\maketitle
\noindent\textbf{Problem ID:} \texttt{H11-187}. \textbf{Status: unresolved; rigorous restricted research attempt.}

\section{Target and restriction}
The target compares feasible capacity and accountability reforms while
keeping appropriation and repression separate from service performance.
The source review calls for distinguishing politicians and bureaucrats
\cite{B}. We solve a two-actor incentive model, establish a conditional
accountability-first safeguard, and show why terminal average constraints
cannot establish pathwise safety. The result is a model-dependent sequencing
criterion, not a universal recommendation or a causal estimate.

A polity has an administrative resource stock $k>0$, oversight parameter
$a>0$ and a fixed bureaucratic bonus scale $w>0$. An elected politician
chooses appropriated resources $x\in[0,k]$ to maximize
\[
 x-\frac a2x^2.                                               \tag{1}
\]
Oversight therefore changes the politician's expected private cost of
appropriation. After observing $x$, an appointed official chooses effort
$e\geq0$ to maximize
\[
 w(1-x/k)e-\frac12e^2.                                       \tag{2}
\]
The fraction of the promised bonus that can be delivered is $1-x/k$.
Citizen service output is $Q=(k-x)e$. A verified capacity vector is
$K=(k,e,Q)$, respectively an administrative resource stock, an effort-based
legal-processing capacity and completed service volume. These are distinct
units. A scalar service objective may use $v(K)=Q$; no claim equates raw
extraction with welfare. Accountability is the independently chosen scale
$A=a$, appropriation harm is $R=x$, and repression is identically zero in
this restricted class. Actual applications need a separate nonzero
repression model.

Resources appropriated by the politician are a distributional harm index;
we do not treat their face value as an additional resource cost after
subtracting them from the government's available stock. The game specifies
behavior and measured outputs, not an interpersonal welfare sum. Fixed
salaries, investment expenditures and monitoring costs are included in a
separate budget. Since equilibrium bonuses never exceed $w^2$ per period,
a budget reserving that amount plus the declared reform costs is feasible.

\section{Separate political and bureaucratic incentives}
\begin{theorem}[Unique equilibrium and comparative statics]\label{game}
For every $k,a,w>0$, the sequential game has a unique subgame-perfect
outcome
\[
 x(k,a)=\min\{k,1/a\},\quad
 e(k,a)=w\left(1-\frac{x(k,a)}k\right),\quad
 Q(k,a)=\frac{w(k-x(k,a))^2}{k}.                              \tag{3}
\]
Appropriation is nondecreasing in $k$ and nonincreasing in $a$.
Service is nondecreasing in both $k$ and $a$. Raising $k$ can strictly
increase appropriation while leaving service at zero.
\end{theorem}
\begin{proof}
The official's strictly concave quadratic has unique maximizer
$e=w(1-x/k)\geq0$. The politician's strictly concave objective is
independent of the subsequent effort and has unrestricted maximizer
$1/a$. Projection onto $[0,k]$ gives the unique $x$ in (3), completing
backward induction and giving $Q$. Monotonicity of the minimum expression
proves the claims for appropriation.

When $k\leq1/a$, output is zero. When $k>1/a$ it equals
$w(k-2/a+1/(a^2k))$, whose derivative in $k$ is
$w(1-1/(a^2k^2))>0$. The formula is continuous at the boundary, so
output is globally nondecreasing in $k$. At fixed $k$, increasing $a$
reduces $x$ and hence increases $(k-x)^2$, proving the other output
monotonicity. Finally choose $0<k_0<k_1\leq1/a$: appropriation rises
from $k_0$ to $k_1$ while output remains zero. This proves the last claim.
\end{proof}
This last observation is an explicit warning about interpreting higher
revenue or administrative stock as effective service capacity. Here the
politician appropriates the entire increment before the official can use
it. More favorable political preferences would change the formula.

\section{Two reforms, the same endpoint, different path safeguards}
Two irreversible reforms raise $k_0$ to $k_1\geq k_0>0$ and $a_0$ to
$a_1\geq a_0>0$. Exactly one reform is implemented in each of two periods.
Both orders have the same undiscounted investment cost $c_k+c_a$ and may
use the same budget $B\geq c_k+c_a+2w^2$ plus fixed administrative costs.
There is no reform interaction or loss of commitment in the transition
technology. Outcomes within each period are (3). Let $r\geq0$ be a
maximum permitted appropriation in every period, including baseline.

\begin{theorem}[Conditional sequencing dominance for safety]\label{sequence}
Suppose the baseline and common endpoint satisfy
\[
 x(k_0,a_0)\leq r,
 \qquad x(k_1,a_1)\leq r.                                    \tag{4}
\]
Then accountability first is feasible for the pathwise appropriation cap.
Capacity first is feasible if and only if
\[
 \min\{k_1,1/a_0\}\leq r.                                   \tag{5}
\]
Both orders have identical terminal capacity vector, accountability and
appropriation. Therefore no criterion involving only those terminal
outcomes can distinguish their interim safeguard feasibility.
\end{theorem}
\begin{proof}
The accountability-first intermediate state is $(k_0,a_1)$.
Theorem~\ref{game} gives
\[x(k_0,a_1)\leq x(k_0,a_0)\leq r.\]
Baseline and final safety hold by (4). This proves all-period feasibility.
The capacity-first intermediate state is $(k_1,a_0)$, where appropriation
is exactly the left side of (5). Baseline and terminal states are already
safe, so (5) is both necessary and sufficient. The final primitive state
$(k_1,a_1)$ is identical for both orders, and the equilibrium is unique.
Hence all its listed terminal outcomes coincide. A terminal-only rule
cannot encode a difference in an intermediate outcome it does not use.
\end{proof}
For example, $k_0=1/5$, $k_1=1$, $a_0=1$, $a_1=4$, $r=3/10$ satisfy
(4): appropriation is initially $1/5$ and finally $1/4$.
Capacity first produces interim appropriation one and fails the cap;
accountability first produces $1/5$ and passes. This is an exact
hypothetical example, not an estimated parameterization. If only the
terminal expected constraints in the canonical problem are imposed,
both orders pass in this example. A pathwise safeguard is a stronger
variant and is explicitly declared as such.

\section{When does a policy intervention identify the two mechanisms?}
\begin{proposition}[Joint interventions and baseline saturation]
Assume $w$ is known and the primitive game is correctly specified.
An experiment that sets $(k,a)$ identifies its mean outcomes by assignment
randomization and consistency; in this deterministic model they equal
(3). Observing only $x=k$ at one stock $k$ restricts $a$ to
$0<a\leq1/k$ and does not identify its value. Observing $0<x<k$ identifies
$a=1/x$. In the saturated baseline case, an intervention increasing stock
to $k'>1/a$ produces an interior $x$ and hence identifies $a$, provided
oversight is unchanged.
\end{proposition}
\begin{proof}
Random assignment makes the observed outcome distribution under each
assigned bundle equal to its potential-outcome distribution; consistency
then identifies its mean. For structural inversion, (3) gives $x=k$
exactly when $k\leq1/a$, equivalently $0<a\leq1/k$. Every such $a$
yields the same appropriation and zero service at that stock. If $x<k$,
the minimum must be $1/a$, giving $a=1/x$. With $k'>1/a$, this second
case applies at the new stock. The invariance condition is essential:
if the stock reform also changes oversight, the new observation identifies
the new $a$, not the baseline one.
\end{proof}
A bundle experiment therefore identifies assigned outcomes more generally
than it identifies a politician channel. The last structural inversion is
valid only inside (1)--(2). It cannot be transported to a game with
unmeasured political preferences, corruption networks or endogenous audit
severity without additional restrictions.

\begin{proposition}[Average safety does not imply individual safety]
For any $r>0$, a population can satisfy $\E R=r$ yet violate $R\leq r$
with probability one half. If instead $0\leq R\leq M$ and
$\E R\leq b$, then for every $r>0$,
$\Pr\{R>r\}\leq\min\{1,b/r\}$; an average restriction by itself need
not make that probability zero.
\end{proposition}
\begin{proof}
Give half the population harm zero and half harm $2r$, which proves the
first statement (and is compatible with a bound $M\geq2r$).
For the second, $R\geq r\Ind\{R>r\}$ pointwise, so taking expectations
gives $r\Pr(R>r)\leq\E R\leq b$. Probability is also at most one.
The first construction demonstrates why this upper bound generally cannot
be replaced by zero.
\end{proof}

\section{Remaining identification and implementation gaps}
The full question permits conflict, regime change, nonstate actors,
uncertain reform implementation and heterogeneous polities. These are
excluded from the solved game. In particular, accountability-first safety
can fail if oversight erodes, if monitoring itself represses citizens or
if delayed capacity makes commitment infeasible. Independent accountability
and harm measurement is required; naming $a$ an accountability index is a
model convention, not a validated empirical scale. The terminal objective
and the stronger interim safeguard must be reported separately.

The source PDF was retrieved and its concluding third research direction,
printed p.~418, was inspected on 9 October 2026. It motivates separation
of political and bureaucratic actors; it does not assert a universal reform
sequence. All results above are complete restricted-model derivations.
The empirical reform rankings and the catalogue's broad target remain
unresolved.
\begin{thebibliography}{9}
\bibitem{B} T. Besley, R. Burgess, A. Khan and G. Xu (2022),
\emph{Bureaucracy and Development}, Annual Review of Economics 14,
397--424. Section 4.2 and conclusion p.~418, third direction.
\url{https://www.robinburgess.com/s/Bureaucracy_and_Development.pdf}.
DOI: \url{https://doi.org/10.1146/annurev-economics-080521-011950}.
\end{thebibliography}

\end{document}
